% 摘要和关键词

\begin{abstract}
{\zihao{-4}\setstretch{1.5}
\par 随着物联网、边缘计算与端侧智能的发展，智能家居系统正由“单点控制”向“多源感知—边缘处理—平台服务—自然交互”的一体化形态演进。在家庭场景中，环境舒适性与空气健康需求要求系统具备连续感知与可视化能力，而访问控制、隐私安全与交互便捷性又对本地智能与可靠通信提出更高要求。围绕上述问题，本文设计并实现了一套端边云协同的智能家居原型系统。
\par 本文的研究内容包括多传感器环境监测、边缘视觉识别、平台化数据服务与语音交互。系统采用分层与模块化架构：底层以STM32作为采集节点，完成温湿度、气压与VOC等环境参数采集，并以结构化串口协议上报至网关；网关层以ESP32-S3为核心，完成数据融合、气体传感器采样、TFT本地显示、WiFi联网与HTTP数据上报，同时基于MQTT实现指令下发与状态回传；视觉侧引入OpenMV实现端侧人脸采集与识别，支持识别结果上报与图像获取；平台侧基于Node.js实现数据接收、用户登录鉴权、历史数据管理与统计接口，并通过Web页面实现多参数曲线展示与人脸系统交互；交互侧实现语音助手子系统，采用“ASR—LLM—TTS”流水线与工具调用机制，将自然语言意图映射为受控接口请求，实现环境信息查询等任务的语音交互。
\par 实现与联调结果表明，该系统能够稳定完成多传感器数据采集、网络上传与网页端可视化展示，控制链路能够支持对人脸系统的采集、识别、状态查询与图像获取等操作。本文的创新点体现在分层协作的端边云架构设计、多协议（UART/HTTP/MQTT）一体化集成、端侧人脸识别与平台联动，以及面向智能家居场景的LLM工具调用式语音交互方案。本文工作可为电子信息工程专业在嵌入式系统设计、通信协议集成、边缘智能与人机交互方向的综合实践提供参考。
}
\end{abstract}

{\zihao{-4}\setstretch{1.5}
\noindent\textbf{关键词：}智能家居；STM32；ESP32-S3；环境监测；端边云协同；人脸识别；语音助手；MQTT；Node.js
}

\vspace{0.5\baselineskip}
\newpage
{\zihao{-4}\setstretch{1.5}
\begin{center}
\textbf{Abstract}
\end{center}
\par With the development of the Internet of Things, edge computing, and on-device intelligence, smart home systems are evolving from single-point control to an integrated paradigm featuring multi-source sensing, edge processing, platform services, and natural interaction. In home scenarios, continuous environmental sensing and visualization are required for comfort and air-quality awareness, while access control, privacy protection, and convenient interaction demand reliable local intelligence and robust communication. To address these needs, this work designs and implements a smart home prototype based on an edge--cloud collaborative architecture.
\par The research focuses on multi-sensor environmental monitoring, edge-side vision recognition, platform-oriented data services, and voice interaction. A layered and modular design is adopted. At the device layer, an STM32 node collects environmental parameters including temperature, humidity, pressure, and VOC, and reports structured data to the gateway via a UART protocol. At the gateway layer, an ESP32-S3 performs data fusion, gas-sensor sampling, TFT local display, Wi-Fi connectivity, and HTTP data uploading, while MQTT is used for command delivery and status feedback. On the vision side, OpenMV enables on-device face collection and recognition with support for recognition reporting and image retrieval. On the platform side, a Node.js server provides data ingestion, user authentication, historical data management, and statistical APIs, together with web pages for multi-parameter visualization and face-system operations. In addition, a voice assistant subsystem is implemented using an ASR--LLM--TTS pipeline with tool calling, which maps natural-language intents to controlled API requests for tasks such as environmental information queries.
\par Implementation and integration results show that the system can stably perform periodic data acquisition, network uploading, and web-based visualization, and that the control chain supports face-system operations including collection, recognition, status querying, and image retrieval. The main innovations include the layered edge--cloud collaborative architecture, integrated multi-protocol communication (UART/HTTP/MQTT), edge-side face recognition with platform linkage, and an LLM tool-calling based voice interaction scheme for smart home scenarios.
\par\noindent\textbf{Key words:} smart home; STM32; ESP32-S3; environmental monitoring; edge--cloud collaboration; face recognition; voice assistant; MQTT; Node.js
}

\newpage
